\chapter{Исследовательский раздел}

\section{Анализ производительности ПО}

В качестве тестовой сцены для анализа времени генерации кадра был выбран чайник университета Юта (Utah Teapot)\cite{utah_teapot}, помещенный в Корнеллскую коробку (Cornell Box)\cite{cornell_box}. Для освещения был использован точечный источник, помещенный на середину потолка.

\noindent Тестовая сцена состоит из:
\begin{itemize}
    \item Корнеллской коробки;
    \item чайника Юта;
    \item точечного источника освещения.
\end{itemize}

\noindent Параметра генерации кадра:
\begin{itemize}
    \item разрешение кадра $1280\times720$ пикселей;
    \item максимальное число переотражений луча = 2.
\end{itemize}

\noindent Параметры системы, на которой будут проводиться замеры времени:
\begin{itemize}
    \item CPU -- 6C/12T @ 3.7GHz;
    \item RAM -- 32GB DDR5;
    \item OS -- Windows 11.
\end{itemize}

Для распараллеливания программы использовался стандарт OpenMP\cite{openmp}. Выбор обусловлен относительной простостой способа примения технологии. В языке C++ эта технология реализуется с помощью директивы \#pragma.

% \#pragma omp parallel for schedule(dynamic)

Для установки числа потоков используется функция \texttt{omp\_set\_num\_threads(int)} из заголовочного файла \texttt{<omp.h>}. Для получения количества логических ядер процессора используется функция \texttt{omp\_get\_max\_threads()}.

\begin{figure}[H]
    \centering
    \includegraphics[width=1\textwidth]{utah-teapot.jpg}
    \caption{Чайники с разным числом полигонов}
    \label{fig:utah-teapot}
\end{figure}

Для исследования времени генерации кадра были выбраны следующие варианты числа потоков:
\begin{itemize}
    \item Однопоточный режим (1);
    \item Количество потоков, равное числу физических ядер процессора (6);
    \item Количество потоков, равное числу логических ядер процессора (12);
    \item Вдвое большее количество потоков по сравнению с числом логических ядер процессора (24).
\end{itemize}

Следующие графики представляют собой зависимость времени генерации изображения от количества граней в полигональной моделе чайника.

\begin{figure}[ht]
    \centering

    % 1 поток
    \begin{subfigure}[b]{0.45\textwidth}
        \centering
        \begin{tikzpicture}
            \begin{axis}[
                    title={1 поток},
                    xlabel={число граней},
                    ylabel={время (с)},
                    grid=major
                ]
                \addplot coordinates {
                        (108,15.351) (532,53.706) (1166,114.694) (2148,199.755)
                    };
            \end{axis}
        \end{tikzpicture}
    \end{subfigure}
    \hfill
    % 6 потоков
    \begin{subfigure}[b]{0.45\textwidth}
        \centering
        \begin{tikzpicture}
            \begin{axis}[
                    title={6 потоков},
                    xlabel={число граней},
                    ylabel={время (с)},
                    grid=major
                ]
                \addplot coordinates {
                        (108,4.505) (532,15.768) (1166,30.743) (2148,55.789)
                    };
            \end{axis}
        \end{tikzpicture}
    \end{subfigure}
    \hfill
    % 12 потоков
    \begin{subfigure}[b]{0.45\textwidth}
        \centering
        \begin{tikzpicture}
            \begin{axis}[
                    title={12 потоков},
                    xlabel={число граней},
                    ylabel={время (с)},
                    grid=major
                ]
                \addplot coordinates {
                        (108,2.818) (532,9.704) (1166,21.069) (2148,38.631)
                    };
            \end{axis}
        \end{tikzpicture}
    \end{subfigure}
    % 24 потока
    \begin{subfigure}[b]{0.45\textwidth}
        \centering
        \begin{tikzpicture}
            \begin{axis}[
                    title={24 потока},
                    xlabel={число граней},
                    ylabel={время (с)},
                    grid=major
                ]
                \addplot coordinates {
                        (108,2.566) (532,9.595) (1166,18.611) (2148,36.063)
                    };
            \end{axis}
        \end{tikzpicture}
    \end{subfigure}

    \caption{Время генерации кадра в зависимости от числа потоков и числа граней}
    \label{fig:all_plots}
\end{figure}

\begin{table}[ht]
    \centering
    \caption{Время генерации кадра в зависимости от числа потоков и числа граней}
    \begin{tabular}{|c|c|c|c|c|}
        \hline
        Число граней & 1 поток & 6 потока & 12 потоков & 24 потоков \\
        \hline
        108          & 15.4    & 4.5      & 2.5        & 2.6        \\
        532          & 53.7    & 15.8     & 8.8        & 9.6        \\
        1166         & 114.7   & 30.7     & 18.8       & 18.6       \\
        2148         & 199.8   & 55.8     & 38.6       & 36.1       \\
        \hline
    \end{tabular}
    \label{tab:frame_times_1dp}
\end{table}


На графике ниже представлена зависимость времени генерации кадра от числа потоков.

\begin{tikzpicture}
    \begin{axis}[
            title={Время генерация кадра в зависимости от числа потоков},
            xlabel={число потоков}, ylabel={время, (с)},
            grid=major
        ]
        \addplot[green, mark=o] coordinates {
                (1,15.351) (2,7.966) (4,4.505) (8,2.818) (12,2.488) (16,2.652) (20,2.515) (24,2.566)
            };

        \addplot[yellow, mark=*] coordinates {
                (1,53.706) (2,27.058) (4,15.768) (8,9.704) (12,8.803) (16,9.455) (20,8.843) (24,9.595)
            };

        \addplot[orange, mark=x] coordinates {
                (1,114.694) (2,57.741) (4,30.743) (8,21.069) (12,18.769) (16,19.316) (20,18.568) (24,18.611)
            };

        \addplot[red, mark=+] coordinates {
                (1,199.755) (2,105.847) (4,55.789) (8,38.631) (12,38.612) (16,33.443) (20,36.925) (24,36.063)
            };

        \legend{108 граней, 532 грани, 1166 граней, 2148 граней}
    \end{axis}
\end{tikzpicture}


\section*{Выводы}
Анализ представленных графиков позволяет сделать вывод, что прирост производительности с увеличением числа потоков носит ограниченный характер. Наиболее эффективным числом потоков является количество логических ядер процессора, при котором достигается максимальная параллельная загрузка. При дальнейшем увеличении числа потоков наблюдается снижение производительности, что связано с дополнительными затратами процессорного времени на переключение контекста между потоками.

% \section{Сравнение производительности на CPU/GPU}

% \section{Сравнение производительности без/с BVH}
% \section{Сравнение изображений, полученных с моделями освещения Ламберта/Фонга/Блинна-Фонга}
% \section{Сравнение изображений c разной гаммо-корреккций/тон-маппингом}

